\iffalse
\section{Challenges and Research Frontiers}
\label{sec:frontiers}

Agentic visual generation is developing along three broad paths: expanding the scope of visual agency, increasing the intelligence available to visual agents, and establishing reliable evaluation. Scope expansion concerns the artifacts, project durations, tools, environments, application consequences, and related operating conditions that a system can handle. Intelligence improvement concerns foundation models, system architectures, reasoning, process training, continual learning, collaboration, deployment, and related capabilities. Evaluation concerns the quality, behavior, cost, and consequences of the complete creation process. The three paths also interact: longer projects tend to expose control failures that short tasks leave hidden, a stronger controller magnifies the cost of weak generation and weak verification, and claims about planning, memory, or self-improvement call for evidence at the process level. The subsections below organize each path around concrete failure modes and the research directions they motivate.

\subsection{How to Expand the Scope of Visual Agency}

% Keep the frontier chapter concise in the contents while retaining detailed headings in the text.
\addtocontents{toc}{\protect\setcounter{tocdepth}{2}}

\subsubsection{Multimodal Creative Workflows}

A multimodal creative workflow produces several related artifacts during one project. Text, images, audio, video, source code, layouts, data, and interactive states all contribute to the final work, and each artifact carries its own representation, editing interface, and failure surface. Producing each artifact is within reach of current systems; a harder problem lies in preserving the relations among them. A character that appears in a storyboard, a narration script, a generated video, and a subtitle track is one identity distributed across four representations, and a change to any one of them propagates to the others through constraints that current systems do not represent explicitly.

Existing systems handle each modality through a separate operation with a limited handoff. The dominant pattern is serial conditioning: a text plan drives an image model, whose output conditions a video model, whose frames are passed to a voice or editor stage. Each handoff re-encodes the artifact into a form that loses part of the upstream structure, and a revision in an early stage forces regeneration of everything downstream. The technical causes are threefold. Identity is re-derived at every generation from a reference image or a textual description, and this is a main reason long-form storytelling still exhibits character drift: facial features, wardrobe, and even body proportions subtly change between scenes because the generative process has no persistent binding to a subject. Constraints such as ``the same room'', ``the same brand palette'', or ``the same data underlying the chart'' are expressed only in the prompt and therefore decay through the chain. The dependency structure among artifacts is implicit in the sequence of calls, so the system has no way to determine which outputs must be refreshed when one input changes.

A useful workflow abstraction treats a project as a graph of linked artifacts with typed dependencies, shared references, and version-stable identities. Under this view, the interesting research problems become concrete: how to factorize a subject into a representation that survives re-rendering across styles, camera angles, and modalities; how to express cross-modal constraints so that a verifier can check them at each step; how to localize an update so that regeneration touches only the artifacts whose dependencies changed, in the way a build system invalidates targets; and how to preview the propagation of a proposed edit before executing it. Progress on these problems would turn a multimodal workflow from a chain of independent calls into an editable artifact in its own right.

\subsubsection{Long-Horizon Creative Projects}

A long-horizon project contains many decisions distributed across turns, scenes, files, and production stages. Later actions depend on earlier requirements, intermediate results, unresolved issues, and user preferences, and the system must carry these forward while the context grows. Short-task success gives limited information about this ability: recent empirical studies of long-horizon agents report a consistent pattern in which per-step success that looks strong on isolated operations drops sharply once tasks require an extended, interdependent action sequence, and the dominant failure categories shift from execution errors toward planning errors, loss of early constraints, and error accumulation through history.

Creative generation inherits these failure modes and adds visual ones. Identity drift, layout drift, obsolete references, and repeated repairs are the visual counterparts of forgetting a constraint or acting on a stale assumption. A system may preserve the conversation while losing the exact artifact version that a decision referred to; it may continue executing a plan after the user has revised a central requirement, because the constraint is still in the context window and no longer attended to during generation. The underlying mechanisms are understood: attention over a long trajectory dilutes early instructions, compressed memory keeps the language while dropping the visual state it referred to, and project-level state remains conflated with dialogue history.

Closing this gap requires project-level state that is independent of the conversation. Such state would connect requirements, artifacts, versions, decisions, and unresolved issues in an explicit record; it would separate stable commitments from provisional choices; and it would expose the dependencies that a proposed revision will touch before it executes. Checkpointing, branching, and rollback---the standard machinery of versioned creative tools---have direct analogues in agent memory, but they need to be lifted from file-level snapshots to decision-level provenance: which observation justified this edit, which requirement it served, and which later decisions depend on it. Equally important is an explicit notion of commitment hygiene: detecting when a later instruction contradicts an earlier one, asking whether the earlier commitment should be revised, and propagating that revision across the affected decisions. Long-form video, multi-page documents, interactive stories, and complex image editing all become tractable in proportion to how much of this bookkeeping the system can carry reliably.

\subsubsection{Open Tool Environments}

An open tool environment contains changing models, applications, code runtimes, browsers, simulators, and data services whose interfaces, costs, permissions, and output formats vary across tasks and over time. Many visual agents are built against a fixed tool library with known interfaces, and the distance between that assumption and the deployed world now limits several domains: GUI agents that perform well on the applications seen during training drop sharply on unseen sites and cross-application workflows, because grounding, action semantics, and layout priors transfer poorly; a version bump in an image model can change aspect-ratio behavior, safety filters, or default styles in ways that silently break downstream stages; and a tool that returns a valid file in the wrong representation forces an extra conversion step that the planner did not budget for.

Three capabilities separate an open-environment agent from a closed-library one. The first is interface acquisition: maintaining capability descriptions for tools, learning their argument and return conventions from documentation and execution feedback, and treating a tool's first use in a project as a probe. The second is cost-aware selection: choosing among tools that satisfy a request by expected quality, latency, financial cost, side effects, and the representation that the next stage requires. The third is graceful degradation: when a preferred tool is unavailable or misbehaves, re-routing to a substitute, falling back to a more conservative operation, or surfacing the failure to the user with an explanation. Each capability has a direct research formulation---tool documentation as a learning signal, tool selection as decision making under uncertainty, and tool failure as a first-class state in the controller---and they remain largely absent from current systems, which still assume a static action space.

Open environments also shift the trust model. Tools carry their own failure modes: a browser carries untrusted content, an IDE executes code, a simulator returns observations that a planner will act on, and an external model can return outputs that violate the project's constraints. An agent that treats every tool output as ground truth inherits whatever the tool got wrong, including adversarial content embedded in retrieved pages or rendered screenshots. Treating tool outputs as evidence to be verified, with provenance and confidence attached, supports the move from closed demonstrations toward practical deployment.

\subsubsection{Embodied and High-Consequence Domains}

In an embodied or high-consequence domain, a generated artifact supports an action, a decision, a physical process, or a regulated activity. The artifact is evaluated through domain constraints and downstream consequences in addition to visual appearance, and the cost of a plausible-looking error scales with the stakes of the setting. Visual agents are entering navigation, physical simulation, engineering, scientific analysis, education, and public communication, and each domain imposes constraints that general-purpose models do not carry: a sim-to-real gap between what a generated scene predicts and how the physical world responds; a requirement that technical drawings, medical illustrations, and safety diagrams satisfy standards that appearance metrics leave unmeasured; and an asymmetry between reversible and irreversible actions that current controllers rarely represent explicitly.

A central technical bottleneck is the gap between generation and prediction. Photorealism and world-modeling are separate properties: empirical studies of physics in generated video show that current systems capture correlations from training data, generalize within the training distribution, and fail when the initial conditions leave it, a pattern consistent with case-based recall of training examples. Treating a generative model as a simulator requires properties that have to be established experimentally: causal, action-conditioned prediction; persistent state over an extended horizon; and physical accuracy on the events that matter for the downstream decision. These are also the properties that distinguish a world model from a video generator, and progress on them transfers directly to planning, since the same model can then be used to roll out candidate actions before committing to one.

High-consequence deployment also requires an explicit autonomy policy. The acceptable level of autonomy depends on the consequence of an error, and a system that applies a single trust threshold across actions treats a low-cost retry and an irreversible release alike. Risk-adjusted autonomy assigns different approval requirements to different classes of action, monitors domain constraints during execution, exposes calibrated uncertainty so that a human reviewer knows when to intervene, and preserves a complete record of the assumptions, references, and checks behind a result. These are control decisions that the architecture has to support natively, and they shape how far a visual agent can move past the demonstration regime.

\subsection{How to Increase the Intelligence of Visual Agents}

\subsubsection{Generative and Multimodal Foundation Models}

Foundation models provide the perceptual and generative capabilities that every higher-level agent decision depends on. Their failure modes are well documented and persist across scales: spatial relations, counting, in-image text, precise local edits, identity preservation under change, rare factual knowledge, and physically correct dynamics remain weak spots even in models that produce photorealistic single samples. A related structural weakness is the mismatch between understanding and generation. A model that interprets an image accurately may have no interface for making an exact change to the underlying representation, so an edit that a human would describe as a one-parameter adjustment becomes a re-generation with collateral change to preserved regions; a model with strong generation may in turn lack the introspection to know which region of its own output violates the instruction. These failures share a root cause: appearance quality is a training signal that does not force the model to internalize the structure that makes an edit local, an identity stable, or a scene physically consistent.

Architecturally, the field is converging on unified models that handle understanding and generation in one backbone, in three broad forms: pure autoregressive, pure diffusion, and hybrid designs that autoregress over semantic tokens while diffusing over pixels. The unification matters for agents because it removes a handoff: a single representation can serve as both the medium of perception and the target of synthesis, which is what makes edits addressable and observations usable as evidence. Unification by itself leaves the capability gaps above in place; it gathers them into one architecture. The sharper formulation is representational: whether the model carries an explicit spatial, temporal, or structural state that an edit can address, or whether each edit must re-sample a global latent wholesale.

Foundation-model progress in an agentic setting calls for a different measurement from a single-shot setting. A more informative question is which downstream agent decisions improve when the foundation model changes under a fixed controller and budget: fewer repair loops for the same edit, fewer misread observations during tool use, fewer violations of preserved regions, and more reliable stopping when the verifier has enough evidence. This keeps the foundation model tied to the behavior the agent is meant to support.

\fi
\section{Challenges and Research Frontiers}
\label{sec:frontiers}

Agentic visual generation is developing along three broad paths: expanding the scope of visual agency, increasing the intelligence available to visual agents, and establishing reliable evaluation. Scope expansion concerns the artifacts, project durations, tools, environments, application consequences, and related operating conditions that a system can handle. Intelligence improvement concerns foundation models, system architectures, reasoning, process training, continual learning, collaboration, deployment, and related capabilities. Evaluation concerns the quality, behavior, cost, and consequences of the complete creation process. The three paths also interact: longer projects tend to expose control failures that short tasks leave hidden, a stronger controller magnifies the cost of weak generation and weak verification, and claims about planning, memory, or self-improvement call for evidence at the process level. The subsections below organize each path around concrete failure modes and the research directions they motivate.

\subsection{How to Expand the Scope of Visual Agency}

% Keep the frontier chapter concise in the contents while retaining detailed headings in the text.
\addtocontents{toc}{\protect\setcounter{tocdepth}{2}}

\subsubsection{Multimodal Creative Workflows}

A multimodal creative workflow produces several related artifacts during one project. Text, images, audio, video, source code, layouts, data, and interactive states all contribute to the final work, and each artifact carries its own representation, editing interface, and failure surface. Producing each artifact is within reach of current systems; a harder problem lies in preserving the relations among them. A character that appears in a storyboard, a narration script, a generated video, and a subtitle track is one identity distributed across four representations, and a change to any one of them propagates to the others through constraints that current systems do not represent explicitly.

Existing systems handle each modality through a separate operation with a limited handoff. The dominant pattern is serial conditioning: a text plan drives an image model, whose output conditions a video model, whose frames are passed to a voice or editor stage. Each handoff re-encodes the artifact into a form that loses part of the upstream structure, and a revision in an early stage forces regeneration of everything downstream. The technical causes are threefold. Identity is re-derived at every generation from a reference image or a textual description, and this is a main reason long-form storytelling still exhibits character drift: facial features, wardrobe, and even body proportions subtly change between scenes because the generative process has no persistent binding to a subject. Constraints such as ``the same room'', ``the same brand palette'', or ``the same data underlying the chart'' are expressed only in the prompt and therefore decay through the chain. The dependency structure among artifacts is implicit in the sequence of calls, so the system has no way to determine which outputs must be refreshed when one input changes.

A useful workflow abstraction treats a project as a graph of linked artifacts with typed dependencies, shared references, and version-stable identities. Under this view, the interesting research problems become concrete: how to factorize a subject into a representation that survives re-rendering across styles, camera angles, and modalities; how to express cross-modal constraints so that a verifier can check them at each step; how to localize an update so that regeneration touches only the artifacts whose dependencies changed, in the way a build system invalidates targets; and how to preview the propagation of a proposed edit before executing it. Progress on these problems would turn a multimodal workflow from a chain of independent calls into an editable artifact in its own right.

\subsubsection{Long-Horizon Creative Projects}

A long-horizon project contains many decisions distributed across turns, scenes, files, and production stages. Later actions depend on earlier requirements, intermediate results, unresolved issues, and user preferences, and the system must carry these forward while the context grows. Short-task success gives limited information about this ability: recent empirical studies of long-horizon agents report a consistent pattern in which per-step success that looks strong on isolated operations drops sharply once tasks require an extended, interdependent action sequence, and the dominant failure categories shift from execution errors toward planning errors, loss of early constraints, and error accumulation through history.

Creative generation inherits these failure modes and adds visual ones. Identity drift, layout drift, obsolete references, and repeated repairs are the visual counterparts of forgetting a constraint or acting on a stale assumption. A system may preserve the conversation while losing the exact artifact version that a decision referred to; it may continue executing a plan after the user has revised a central requirement, because the constraint is still in the context window and no longer attended to during generation. The underlying mechanisms are understood: attention over a long trajectory dilutes early instructions, compressed memory keeps the language while dropping the visual state it referred to, and project-level state remains conflated with dialogue history.

Closing this gap requires project-level state that is independent of the conversation. Such state would connect requirements, artifacts, versions, decisions, and unresolved issues in an explicit record; it would separate stable commitments from provisional choices; and it would expose the dependencies that a proposed revision will touch before it executes. Checkpointing, branching, and rollback---the standard machinery of versioned creative tools---have direct analogues in agent memory, but they need to be lifted from file-level snapshots to decision-level provenance: which observation justified this edit, which requirement it served, and which later decisions depend on it. Equally important is an explicit notion of commitment hygiene: detecting when a later instruction contradicts an earlier one, asking whether the earlier commitment should be revised, and propagating that revision across the affected decisions. Long-form video, multi-page documents, interactive stories, and complex image editing all become tractable in proportion to how much of this bookkeeping the system can carry reliably.

\subsubsection{Open Tool Environments}

An open tool environment contains changing models, applications, code runtimes, browsers, simulators, and data services whose interfaces, costs, permissions, and output formats vary across tasks and over time. Many visual agents are built against a fixed tool library with known interfaces, and the distance between that assumption and the deployed world now limits several domains: GUI agents that perform well on the applications seen during training drop sharply on unseen sites and cross-application workflows, because grounding, action semantics, and layout priors transfer poorly; a version bump in an image model can change aspect-ratio behavior, safety filters, or default styles in ways that silently break downstream stages; and a tool that returns a valid file in the wrong representation forces an extra conversion step that the planner did not budget for.

Three capabilities separate an open-environment agent from a closed-library one. The first is interface acquisition: maintaining capability descriptions for tools, learning their argument and return conventions from documentation and execution feedback, and treating a tool's first use in a project as a probe. The second is cost-aware selection: choosing among tools that satisfy a request by expected quality, latency, financial cost, side effects, and the representation that the next stage requires. The third is graceful degradation: when a preferred tool is unavailable or misbehaves, re-routing to a substitute, falling back to a more conservative operation, or surfacing the failure to the user with an explanation. Each capability has a direct research formulation---tool documentation as a learning signal, tool selection as decision making under uncertainty, and tool failure as a first-class state in the controller---and they remain largely absent from current systems, which still assume a static action space.

Open environments also shift the trust model. Tools carry their own failure modes: a browser carries untrusted content, an IDE executes code, a simulator returns observations that a planner will act on, and an external model can return outputs that violate the project's constraints. An agent that treats every tool output as ground truth inherits whatever the tool got wrong, including adversarial content embedded in retrieved pages or rendered screenshots. Treating tool outputs as evidence to be verified, with provenance and confidence attached, supports the move from closed demonstrations toward practical deployment.

\subsubsection{Embodied and High-Consequence Domains}

In an embodied or high-consequence domain, a generated artifact supports an action, a decision, a physical process, or a regulated activity. The artifact is evaluated through domain constraints and downstream consequences in addition to visual appearance, and the cost of a plausible-looking error scales with the stakes of the setting. Visual agents are entering navigation, physical simulation, engineering, scientific analysis, education, and public communication, and each domain imposes constraints that general-purpose models do not carry: a sim-to-real gap between what a generated scene predicts and how the physical world responds; a requirement that technical drawings, medical illustrations, and safety diagrams satisfy standards that appearance metrics leave unmeasured; and an asymmetry between reversible and irreversible actions that current controllers rarely represent explicitly.

A central technical bottleneck is the gap between generation and prediction. Photorealism and world-modeling are separate properties: empirical studies of physics in generated video show that current systems capture correlations from training data, generalize within the training distribution, and fail when the initial conditions leave it, a pattern consistent with case-based recall of training examples. Treating a generative model as a simulator requires properties that have to be established experimentally: causal, action-conditioned prediction; persistent state over an extended horizon; and physical accuracy on the events that matter for the downstream decision. These are also the properties that distinguish a world model from a video generator, and progress on them transfers directly to planning, since the same model can then be used to roll out candidate actions before committing to one.

High-consequence deployment also requires an explicit autonomy policy. The acceptable level of autonomy depends on the consequence of an error, and a system that applies a single trust threshold across actions treats a low-cost retry and an irreversible release alike. Risk-adjusted autonomy assigns different approval requirements to different classes of action, monitors domain constraints during execution, exposes calibrated uncertainty so that a human reviewer knows when to intervene, and preserves a complete record of the assumptions, references, and checks behind a result. These are control decisions that the architecture has to support natively, and they shape how far a visual agent can move past the demonstration regime.

\subsection{How to Increase the Intelligence of Visual Agents}

\subsubsection{Generative and Multimodal Foundation Models}

Foundation models provide the perceptual and generative capabilities that every higher-level agent decision depends on. Their failure modes are well documented and persist across scales: spatial relations, counting, in-image text, precise local edits, identity preservation under change, rare factual knowledge, and physically correct dynamics remain weak spots even in models that produce photorealistic single samples. A related structural weakness is the mismatch between understanding and generation. A model that interprets an image accurately may have no interface for making an exact change to the underlying representation, so an edit that a human would describe as a one-parameter adjustment becomes a re-generation with collateral change to preserved regions; a model with strong generation may in turn lack the introspection to know which region of its own output violates the instruction. These failures share a root cause: appearance quality is a training signal that does not force the model to internalize the structure that makes an edit local, an identity stable, or a scene physically consistent.

Architecturally, the field is converging on unified models that handle understanding and generation in one backbone, in three broad forms: pure autoregressive, pure diffusion, and hybrid designs that autoregress over semantic tokens while diffusing over pixels. The unification matters for agents because it removes a handoff: a single representation can serve as both the medium of perception and the target of synthesis, which is what makes edits addressable and observations usable as evidence. Unification by itself leaves the capability gaps above in place; it gathers them into one architecture. The sharper formulation is representational: whether the model carries an explicit spatial, temporal, or structural state that an edit can address, or whether each edit must re-sample a global latent wholesale.

Foundation-model progress in an agentic setting calls for a different measurement from a single-shot setting. A more informative question is which downstream agent decisions improve when the foundation model changes under a fixed controller and budget: fewer repair loops for the same edit, fewer misread observations during verification, fewer collapsed candidates during search. A weaker generator taxes the controller with repeated retries and inflates the planning burden; a weaker interpreter corrupts the agent's task representation before any tool is invoked. Treating foundation-model improvement as one factor in a controlled agentic evaluation connects model progress to system progress.

\subsubsection{Agent Architecture and Control}

Agent architecture determines how task information, visual state, plans, tools, observations, and actions are organized, and the practical design space spans external orchestration over frozen models, native visual agents, hierarchical workflows, role-specialized teams, executable programs, and hybrids of these forms. The recurring structural failure across all of them is that a collection of components still behaves as a fixed pipeline: information is lost at each interface, a specialist optimizes a local objective that conflicts with the project, a coordinator repeats the same action after different failures, and the boundary between planning, execution, and correction is diffuse enough that responsibility for a failure is hard to localize.

The architectural question that matters is where decisions live and what state they share. Complex visual tasks require decisions at several levels---project organization, operation selection, local editing, verification, delivery---and these levels have different time constants, different information requirements, and different costs of being wrong. An architecture that funnels all of them through a single context window couples them unnecessarily; an architecture that isolates them into uncommunicating specialists loses the cross-level constraints that make the result coherent. The useful middle ground makes state transitions explicit: a shared project state that all roles read from and write to through typed interfaces, decision rights that specify which level can commit which kind of change, and evidence requirements that specify what an observation must contain to justify a transition. This resembles how mature software systems manage concurrency more than how current agents manage turns.

Two further directions follow from treating the architecture as a control problem. Adaptive depth allocates reasoning and search according to uncertainty and consequence, so routine edits take the fast path and ambiguous or risky decisions trigger deeper deliberation---the principle behind test-time compute scaling in reasoning models, applied here to the action loop. Reconfigurable roles treat agents and tools as replaceable components with declared preconditions and effects, so that a role can be created, merged, suspended, or substituted as the task demands, and a failure in one component can be isolated. Both require the architecture to know, at each step, what it is doing and why---a property that current systems approximate through prompts and rarely enforce through structure.

\subsubsection{Visual Reasoning, World Knowledge, and Causal Understanding}

Visual reasoning connects language, visual evidence, structured relations, external knowledge, and the space of possible actions; causal understanding adds a prediction of how an artifact or environment will change under an intervention. Current systems produce plausible appearances with incorrect relations, dimensions, event orders, and physical consequences, and the failure has practical weight: it reflects the gap between correlation and mechanism, in that a model that has seen falling objects has not necessarily internalized gravity. The empirical evidence is consistent across settings---generative models trained on physical dynamics generalize inside the training distribution and fail on out-of-distribution initial conditions, indicating case-based internalization of training examples.

For an agent, this gap determines whether the system can plan. Rolling out a candidate edit internally, predicting what a scene will look like after a motion, or forecasting how a diagram will break when a value changes all require an action-conditioned predictive model, and the properties such a model needs---causality, interactivity, persistent state over long horizons, physical accuracy---are largely the properties that current video generators lack. The research trajectory on world models makes the requirements explicit and the trade-offs measurable: causal autoregressive structure enables interaction, persistent memory enables coherence across an extended rollout, and physical fidelity trades against real-time responsiveness. Different applications weight these differently, but an agent that intends to plan in the world needs at least the first two.

Closing the gap between appearance and mechanism calls for representations that make structure explicit. Scene graphs, layouts, executable programs, geometric and simulation states, retrieval-augmented factual grounding, and explicit uncertainty each shift part of the reasoning burden out of the latent space into a form that can be inspected and verified. Training can reinforce the same shift: counterfactual supervision that varies an action while holding the initial state fixed teaches the model what depends on what, and evaluation that scores the relation between an output and the process that produced it rewards the same structure. The end state is a model for which ``the result is correct'' and ``the result follows from the input the way the task intended'' coincide, which is what the term visual reasoning should mean in an agentic setting.

\subsubsection{Agentic Reinforcement Learning and Process Training}

Process training teaches an agent how to make a sequence of decisions: which tool to call, how to construct a prompt, when to inspect, when to revise, and when to stop. Reinforcement learning connects those decisions to outcomes observed later in the trajectory. The central difficulty is credit assignment: a final-artifact score provides weak information about intermediate choices, a useful repair several steps earlier receives little signal, and a lucky sample can receive more than a deliberate one. Learned evaluators introduce a second difficulty---specification gaming---whose documented forms in agentic settings include exploiting the verifier, optimizing a proxy metric while damaging user intent, preserving an already-correct region that a metric happens to reward, and, in the sharpest cases, modifying the evaluation itself. The empirical pattern is consistent: the stronger the optimizer and the weaker the reward, the more reliably the gap between the metric and the intent is exploited.

Process rewards are the standard response and remain an open problem. Step-level supervision requires annotations or evaluators that can judge an intermediate decision, which is hard for visual work: whether a tool call was correct depends on information that becomes available only later, and whether an intermediate image is on track depends on the intended final state. Verifier-guided training and test-time search improve the signal but concentrate effort on the verifier, which then becomes both the bottleneck and the attack surface. A robust training recipe therefore combines signals with different failure modes---outcome rewards, process rewards, preference comparisons, environment feedback, counterfactual branches, failure injection---so that no single evaluator's shortcuts dominate the learned policy.

Reward design for visual agents also has to encode the properties that make agentic creation distinct from single-shot generation: constraint satisfaction, collateral change to preserved regions, resource expenditure, willingness to ask for help, and the cost of an irreversible action taken under uncertainty. These belong among the primary objectives of an agent operating over a long trajectory, and a reward that omits them teaches a policy that succeeds on the benchmark and struggles in deployment. The research frontier lies in reward and verifier design that survives optimization, and in evaluation that can tell a policy that learned the task from one that learned the test.

\subsubsection{Memory, Continual Learning, and Self-Evolution}

Memory makes information from earlier operations available to later decisions; continual learning changes how the agent behaves across tasks; self-evolution creates or updates reusable strategies, skills, tools, and routing rules. The three are often discussed together and fail for different reasons. Memory in current systems is largely a retrieval problem over dialogue and artifacts, and its documented failure is state tracking: as facts, constraints, and decisions are revised over a long interaction, an agent that retrieves by similarity returns the superseded state as readily as the current one, because similarity to the query is a poor signal for what is true now. Long-horizon failure studies place memory limitations and catastrophic forgetting among the dominant categories, with a specific mechanism---an early constraint remains in the context window and is no longer attended to during generation---that persists even as context windows grow.

Continual learning across visual projects has an additional structural requirement: the experience has to be represented at a useful level of abstraction. A stored trajectory tied to one tool's parameter names transfers poorly; a stored strategy with explicit preconditions transfers well. The reusable units of visual work are skills, repair patterns, tool compositions, and constraint templates, and the interesting questions are how to extract them from successful trajectories, how to attach the conditions under which they apply, and how to detect when an accumulated skill has become stale because the underlying tool or model changed. A skill library that lacks versioning reproduces yesterday's workarounds against today's interfaces and turns into a liability over time.

Self-evolution sharpens both problems because an update that goes unmeasured becomes a hidden change of policy. A benchmark gain that comes from extra retrieval or sampling while the underlying behavior stays fixed reflects test-time adaptation; a system that absorbs an error into its memory reproduces it across future tasks. The safeguards are the ones established in continual learning---held-out evaluation, forgetting and contamination checks, versioned updates with rollback---applied to behavioral change, which is the unit that matters when the learned artifact is a skill library or a router. The practical payoff for visual agents is concrete: fewer re-derived tool paths for repeated project types, personalization to a user's stylistic commitments, and transfer of a repair pattern from one domain to another, all of which require memory that is structured, evaluated, and reversible.

\subsubsection{Multi-Agent Collaboration and Human Co-Creation}

Multi-agent collaboration divides interpretation, generation, critique, retrieval, editing, and delivery across several roles; human co-creation adds user goals, preferences, contextual knowledge, evaluation, and authorship to the same process. Multi-agent designs promise specialization and breadth; their recurring cost is coordination, which shows up as duplicated work, disagreement about the current state, optimization of incompatible local objectives, communication that rivals the compute budget of the generator itself, and error propagation through handoffs. Failure analyses of multi-agent systems repeatedly identify the same mechanism: without a shared, authoritative state, each role maintains its own picture of the project, and the divergence surfaces in the final artifact.

Human co-creation inherits the coordination problem and adds an asymmetry of visibility. A user whose visibility ends at the final output has no way to tell whether an important decision was made at the specification level, where their preference would have changed it, or in an intermediate step, where changing it now is expensive. The consequence is either micro-supervision, in which the user checks every low-level operation because the system does not escalate, or silent divergence, in which the system produces something polished but off-target. Both failure modes trace back to architecture: the current designs lack an explicit policy for when the system asks, when it proposes alternatives, when it defers, and when it proceeds---an adaptive initiative that tracks the risk and reversibility of the current decision.

The design levers are correspondingly structural. Shared state with typed ownership makes the current project state authoritative and reduces cross-role divergence. Explicit communication protocols bound the coordination cost. Role-specific evidence assigns each specialist the observations it is entitled to act on. Conflict resolution specifies what happens when two roles disagree. User-selectable levels of intervention let a person tune the granularity at which they participate, from high-level direction to local edits, with the system surfacing trade-offs alongside outcomes. Evaluation of such systems has to measure coordination cost, user workload, and authorship clarity alongside task quality, because a multi-agent architecture that improves output while making the process opaque has traded a visible problem for an invisible one.

\subsubsection{Efficiency and Scalable Deployment}

Efficient deployment concerns the quality, latency, compute, energy, and human time required to operate an agentic visual system. The cost structure of an agent differs from that of a single model call in ways that shape deployment: model calls, tool calls, retrieval, simulation, multi-agent communication, candidate generation, and review compound over a trajectory, and in visual work the candidates are large---video, 3D scenes, layered documents---so the marginal cost of one more try is far higher than in text. Multi-agent communication and visual verification can rival or exceed the cost of the generator itself, and a workflow that performs acceptably in a demonstration can become impractical for long videos, dense 3D scenes, or many concurrent users.

The principled response is to allocate compute according to uncertainty and consequence. Test-time scaling in visual generation has matured into an explicit trade-off between sample count, verifier quality, and final performance, and the empirical lesson parallels the one from reasoning models: verification-guided scaling is more sample-efficient than blind sampling, and a calibrated stopping rule recovers most of the quality at a fraction of the cost. The same logic extends to the action loop---a useful agent estimates the expected value of another action, selects the depth of verification, reuses cached results, and chooses between fast and slow paths---so that effort concentrates on the decisions where it changes the outcome. Fixed retry budgets remain the default in many systems and are the least developed part of this design space.

Deployment also constrains which architectural and learning choices are viable, and the constraint runs in both directions. A causally structured, few-step generation process makes interactive visual systems feasible and trades against the physical fidelity that high-stakes applications need, and different deployments weight the two differently. Serving infrastructure that treats generation, verification, and memory as independently optimized stages leaves the joint optimization unexploited; reporting quality in isolation hides the quality--cost and quality--latency trade-off that deployment actually faces. Energy, human review time, and the cost of rejected candidates belong in the same report as model expenditure, because for an agentic system they are of comparable magnitude and are currently the least visible part of the bill.

\subsection{How to Establish Reliable Evaluation}

\subsubsection{Beyond Final-Artifact Quality}

Final-artifact evaluation measures the appearance, semantic alignment, or task quality of the output; agent evaluation also measures how the system interpreted the task, chose actions, used evidence, handled resources, and involved people. The reason is attribution: a final score leaves open whether improvement came from the foundation model, the controller, extra samples, a human correction, or a lucky branch, and an agentic system that claims planning, reflection, memory, or recovery is making a claim about process behavior that a last-frame metric leaves unobserved. A fixed pipeline that retries the same call can earn the same score as a system that adapts its behavior, which makes the score uninformative about the claim under test.

Process evaluation has to be designed with intent. Which decisions the system faced, what evidence was available at each, what action was taken, what the alternative actions were, and what the action changed are the units of agentic behavior, and a report that separates artifact quality, constraint satisfaction, state retention, action selection, diagnosis accuracy, recovery, stopping, resource use, and human intervention can localize an improvement to the component that produced it. Attribution at this granularity turns an evaluation into a diagnostic instrument: it tells a researcher whether the next unit of effort should go into the foundation model, the controller, the memory, or the verifier. Baselines play a corresponding role---holding the generator, tool set, sampling budget, and human access fixed while varying the controller or learning method is what makes an agentic claim interpretable.

Practice in the field has started to move in this direction, with trajectory-grounded automatic judges used to attribute failures to categories such as planning error, memory limitation, or constraint loss, and with the composition of failure types shifting measurably as task horizon grows. Extending this practice to visual agents requires solving the problems specific to visual trajectories: observations that take the form of images, artifact versions that must be tracked across edits, and intermediate states that lack a compact symbolic form. A process evaluation for visual agents is therefore also a data-engineering problem---what to record, at what fidelity, and how to make a long multimodal trajectory comparable across systems---and the evaluation claims below draw on its solution.

\subsubsection{Process, Causal, and Counterfactual Evaluation}

Process evaluation records the sequence of states, observations, actions, results, and decisions. Causal evaluation asks whether an action produced the observed improvement. Counterfactual evaluation changes an input, an observation, or an available action and tests whether the behavior responds. The three are layered, and current practice sits mostly below the first: reported trajectories routinely omit the observation used at a decision point, the version of a tool, the reason for a retry, or the state before a correction, which leaves open whether feedback changed the action or whether the system would have produced the same output on a fixed path. The defining property of an agentic loop lies in the connection between evidence and action, and a benchmark that leaves this connection unexposed has limited power to separate adaptive control from rehearsed execution.

The concrete design instrument is perturbation. Injecting localized defects, unavailable tools, changed references, invalid programs, conflicting constraints, or misleading observations measures what a system does when the world departs from its expectation, which is the situation an agent is built for. The measurable responses are diagnostic: whether the defect is detected, whether the diagnosis identifies the right cause, whether the repair is local or destructive, whether preserved constraints survive, how long recovery takes, and whether the system stops or escalates appropriately. Comparing the action selected under alternative evidence isolates the causal contribution of the observation to the decision and provides direct evidence on whether the loop is closed.

For visual agents, perturbation has an additional dimension because observations are high-dimensional and generative. A misleading screenshot, an ambiguous reference image, a partially corrupted intermediate artifact, or an adversarial instruction embedded in retrieved content are all realistic observations, and clean-input evaluation leaves the agentic behavior itself largely untested. Benchmarks that treat visual perturbation as a first-class axis, alongside tool and task perturbation, are the corresponding open need, and the payoff extends past robustness: the same machinery that measures recovery from an injected defect measures whether a claimed reflection or self-correction mechanism actually modifies behavior, which is the claim most often made and least often tested in current systems.

\subsubsection{Long-Horizon and Cross-Task Evaluation}

Long-horizon evaluation tests a system over many related decisions; cross-task evaluation tests whether information from earlier tasks affects later ones under new prompts, artifacts, tools, users, or domains. Both are under-served for the same reason: they require evaluation infrastructure that maintains state across episodes, while most benchmarks are collections of independent single-step instances. The empirical consequence is well characterized in adjacent domains---short episodes conceal state drift, memory errors, repeated failures, and stopping problems, and performance degrades as the number of interdependent decisions grows, with planning-related and memory-related failures dominating the composition of errors at longer horizons. Visual creation adds drift modes of its own---identity, layout, style, and reference drift---that accumulate quietly until the artifact is visibly inconsistent with its own earlier state.

Cross-task evaluation carries a sharper trap because self-improvement claims are easy to state and hard to evidence. An update derived from a set of tasks and evaluated on the same distribution shows in-distribution adaptation; evidence for transfer requires held-out tasks. A memory that retrieves the most similar past case can improve a benchmark through lookup while leaving the policy unchanged. Meaningful evidence requires separating formation tasks, development tasks, and held-out transfer tasks, and introducing shifts that the update could not have anticipated: changed requirements, new tools, new visual domains, delayed references. The measurable quantities are then the ones continual learning has standardized---forgetting of prior competence, contamination of the evaluation, persistence of the update, and the relation between a specific stored experience and a specific later decision.

Establishing this evaluation for AVG bears directly on the L5 claims that the autonomy scale in this survey defines. Behavior that changes under experience is difficult to validate with a static test set; the test set has to evolve with the system or be partitioned against it. Long-horizon and cross-task benchmarks for visual agents therefore carry significance well past an expanded evaluation suite: they bear on whether the upper end of the autonomy scale describes current systems or anticipates them.

\subsubsection{Multidimensional Visual Correctness}

Visual correctness spans several partially independent axes, and these axes can conflict. Appearance, semantic alignment, spatial relations, geometry, topology, text, data, identity, temporal ordering, physical plausibility, executable behavior, and edit preservation can each fail on its own, and improving one can degrade another---a more attractive render can break a dimension, a stricter preservation can block a needed repair. A single score compresses these failures and, more consequentially, obscures which requirement failed, which is the information an agent uses to plan a repair. A visually convincing result with an incorrect value, a missing object, an impossible shadow, or a changed identity is precisely the failure mode that appearance-weighted evaluation rewards and that domain users cannot accept.

Different artifacts require different evidence sources. An image, a CAD model, a chart, a page, and an interactive interface each expose correctness through a different channel---rendering, geometric validation, data re-execution, parsing, and interactive testing respectively---and a single visual-similarity metric measures appearance while leaving the other channels largely unexamined. The corresponding evaluation architecture represents requirements as typed claims and associates each claim with the check that can adjudicate it: learned judgments where appearance is the claim, symbolic and programmatic checks where structure or data is, execution and simulation where behavior is, and human review where acceptance is. Preserving per-claim outcomes and confidence alongside any aggregate keeps the diagnostic information that a single number discards.

For an agent, multidimensional correctness has a further role: it is the interface between evaluation and control. A verifier that returns a typed claim---what failed, where, and with what confidence---is actionable, whereas a scalar score forces the controller to guess what to change. Closing the loop between multidimensional evaluation and repair is one of the higher-leverage connections in the field: the same typed-claim representation serves the verifier that guides sampling, the diagnostic that localizes a defect, and the report that tells a user what to trust, so progress here compounds across the stack.

\subsubsection{Human, Cost, Safety, and Provenance Evaluation}

Human, cost, safety, and provenance evaluation measures the conditions under which a visual agent can be trusted and used: user understanding, intervention timing, workload, resource expenditure, untrusted inputs, tool side effects, privacy, authorship, and the history of an artifact. Current reporting under-measures these factors. User studies report preference without control or workload; system reports omit human corrections and the cost of rejected candidates; provenance is lost through editing, delegation, format conversion, and rollback. The result is an evaluation regime in which the cost of producing a deliverable and the conditions under which a person could trust it both stay out of view.

The safety surface of a visual agent is wider than that of a text agent because the input channel is wider. Visual and audio content can carry instructions that the user did not authorize---embedded in retrieved images, rendered pages, documents, or screenshots---and multimodal injection is a documented, reproducible attack against the vision-language backbone that most visual agents sit on. An agent that treats each observation as trusted data can be steered by a single injected image; defenses therefore have to be built into the architecture, separating sources of authority from sources of content, and evaluated under adversarial testing. Provenance and content-credential machinery gives the corresponding record on the output side---what was generated, from what sources, edited by whom---and matters for an agent in particular because its artifacts accumulate edits across a trajectory and the provenance record has to survive that accumulation.

The unifying frame is calibrated trust: what matters is the relationship between system ability, human oversight, and consequence, and a quality score gives limited evidence about it. The measurable quantities are the ones that make that relationship visible: reporting model, tool, retrieval, simulation, communication, and human costs as separate line items; testing user understanding and intervention timing alongside preference; probing permission boundaries, prompt injection, and reversible action coverage as first-class cases; and connecting sources, transformations, responsible actors, and approvals across the full trajectory. These measures bear directly on whether a visual agent can be deployed where it matters, and they are currently the least standardized part of the evaluation stack.

\subsection{More Questions About Agentic Visual Generation}

\subsubsection{What Does Intelligence Mean in Visual Creation?}

Intelligence in visual creation includes understanding, representation, generation, action selection, adaptation, exploration, collaboration, uncertainty management, and responsibility, and the relative value of these capacities depends on the task. A system can be strong in appearance and weak in structural reasoning, strong in local editing and weak in project continuity, strong in exploration and weak in user control, and a single ranking obscures these differences. The same divergence appears across applications: personal artwork, scientific figures, engineering designs, educational media, and high-consequence applications weight creativity, accuracy, speed, autonomy, and oversight differently, so systems optimized for different profiles occupy different points in a capability space.

The practical implication is that intelligence claims are best stated as profiles under explicit conditions: which capacity improved, under what task and budget, at what cost to other capacities, and with what effect on the user and the artifact. Profile reporting makes comparison meaningful across settings that value different things, and it raises the question that aggregate scores set aside: intelligence for whom, and for what. For AVG in particular, where the artifact carries intent and the process carries control, an intelligence claim needs an account of what the system understood, what it changed, and what it preserved before it can be evaluated.

\subsubsection{Authorship, Responsibility, and Human Control}

Agentic creation distributes creative decisions across users, agents, foundation models, tools, data sources, and evaluators, and the distribution raises two questions that a quality metric leaves aside. Authorship concerns contribution and creative direction---who shaped the intent, who selected among alternatives, whose style the result carries. Responsibility concerns who can approve, explain, change, or reverse a consequential decision, and it becomes acute when an agent's output affects external files, public content, collaborative projects, or regulated activity. Multiple models and tools contribute to a single artifact in current systems, and their individual contributions and permissions are difficult to reconstruct after the fact, which makes both questions hard to answer even when the process was well intentioned.

The design goal is meaningful control at the points where values enter the artifact. Editable intermediate states let a user intervene while intent is still structural, before it is baked into pixels. Approval checkpoints attach to consequential and irreversible decisions. Contribution records and provenance connect an output to the sources, models, and decisions that produced it. Permission boundaries and reversible actions define what the system may do unilaterally and what it must escalate. Adaptive initiative ties the level of autonomy to the risk of the decision. Together these mechanisms maintain authorship and responsibility during operation, at the stages where they can still be influenced.

\subsubsection{From Individual Models to Creative Ecosystems}

A creative ecosystem combines foundation models, agent policies, tools, data, memory, evaluation services, interfaces, human workflows, and deployment infrastructure, and the complete system determines what a user can create and how reliably a result can be revised and delivered. Ecosystem-level fragility is already visible: a workflow can depend on a particular model version, external API, evaluator, data source, or application interface, and a change in one layer alters the behavior of the whole pipeline while every local component still works. Communication overhead, licensing, dependency drift, and attribution become the binding constraints as the ecosystem grows, and single-model improvements leave them in place.

Ecosystem-level design is a research object in its own right. Model and tool contracts specify behavior at interfaces so that one component can be replaced without re-deriving the workflow. Version dependencies and provenance records make the behavior of a deployed system reproducible and auditable. Evaluator behavior specified as part of the contract keeps a silent verifier change from invalidating a trained policy. Permission boundaries and resource budgets travel with the workflow. Component substitution that preserves the meaning of the artifact and the trajectory record extends to the ecosystem level the locality that editing seeks at the artifact level. Visual creation is moving toward reusable workflows and shared services, and progress in any single model or tool compounds across many applications in proportion to the stability of these interfaces, provenance mechanisms, and evaluation protocols.

The frontier agenda across all three paths therefore concerns a changing unit of visual intelligence: a single result, a multimodal workflow, a sustained creation process, a system that learns from its history, and a responsible participant in human creative work. Each stage requires evidence about what the system understood, what it changed, what it learned, what it cost, and how people can control its consequences, and each stage adds to the obligations of the previous one.

\addtocontents{toc}{\protect\setcounter{tocdepth}{3}}
